% toy-paper.tex — fixture manuscript for /audit-reproducibility, /verify-claims,
% and /review-paper testing (Workstream B-II). The answer key for the planted
% claims lives in tests/TESTING.md, NOT here — auditors must not be tipped off.
%% source: tests/fixtures/toy-panel/expected_estimates.json
\documentclass[12pt]{article}
\usepackage{amsmath}
\usepackage{natbib}
\usepackage[margin=1in]{geometry}

\title{Treatment Effects in the Toy County Panel:\\ A Fixture Manuscript}
\author{Workflow Test Bed}
\date{June 2026}

\begin{document}
\maketitle

\begin{abstract}
We estimate the effect of the 2015 program on county outcomes using a
balanced county-by-year panel. Two-way fixed effects estimates indicate a
treatment effect of 1.94 units, consistent with program projections.
This manuscript exists to be audited, not believed.
\end{abstract}

\section{Introduction}

Difference-in-differences designs with a single adoption date avoid the
negative-weighting problems that arise under staggered timing
\citep{goodmanbacon2021}. Prior work on placebo program evaluation reports
that county-level interventions of this scale typically yield effects
between 1.5 and 2.5 units \citep{smith2019placebo}.

\section{Data and Design}

The sample is a balanced panel of $N = 200$ county-year observations:
20 counties observed annually from 2010 to 2019. Counties 11--20 adopt the
program in 2015; counties 1--10 never adopt. The estimating equation is
\begin{equation}
  y_{it} = \alpha_i + \gamma_t + \tau \, D_{it} + \varepsilon_{it},
\end{equation}
with standard errors clustered at the county level (20 clusters).

\section{Results}

\subsection{Main estimate}

The two-way fixed effects estimate of the treatment effect is
$\hat{\tau} = 1.71$ (cluster-robust SE $= 0.37$), statistically
distinguishable from zero at the one percent level.

\subsection{Robustness}

Excluding the 2015 transition year --- when treatment exposure is
partial --- the estimated effect is $\hat{\tau} = 1.62$ ($N = 180$),
close to the main estimate.

\section{Conclusion}

The program raised county outcomes by approximately 1.7 units. The
magnitude is in line with the placebo-studies benchmark range documented
by \citet{smith2019placebo}.

\bibliographystyle{plainnat}
\bibliography{toy-paper}

\end{document}
